% Soluzioni dell'esercitazione TEEA 3MAT — Argomento 2: partitori, potenza ed energia
% Per rigenerare il PDF (dalla cartella strumenti, due volte):
%   pdflatex -output-directory=../../../materiali -jobname=soluzioni-esercitazione-partitori-potenza-energia-3MAT soluzioni.tex
% ---------------------------------------------------------------------------
%  Stile comune di esercitazione e soluzioni (pdfLaTeX).
%  Viene incluso da esercitazione.tex e soluzioni.tex con \input.
% ---------------------------------------------------------------------------
\documentclass[11pt,a4paper]{article}
\usepackage[a4paper, top=14mm, bottom=18mm, left=16mm, right=16mm, footskip=9mm]{geometry}
\usepackage[utf8]{inputenc}
\usepackage[T1]{fontenc}
\usepackage{helvet}
\renewcommand{\familydefault}{\sfdefault}
\usepackage{amsmath, amssymb, graphicx, xcolor, fancyhdr, lastpage, tabularx, booktabs, array, enumitem}
\usepackage[most]{tcolorbox}
\usepackage[siunitx, american resistors, american currents]{circuitikz}
\ctikzset{bipoles/thickness=1, current/distance=0.35}
\usetikzlibrary{calc, arrows.meta}

\definecolor{titolo}{HTML}{2A7FA8}
\definecolor{grigio}{HTML}{E9ECEF}
\definecolor{tenue}{HTML}{555555}
\setlength{\parindent}{0pt}
\setlength{\parskip}{3pt}

\newcommand{\U}[1]{\,\mathrm{#1}}
\newcommand{\virg}{{,}}
\newcommand{\riga}[1][28mm]{\underline{\hspace{#1}}}

% intestazione come nelle esercitazioni della scuola
\newcommand{\intestazione}[2]{%
  \begin{minipage}[c]{20mm}\includegraphics[width=18mm]{logo-isis.png}\end{minipage}%
  \begin{minipage}[c]{\dimexpr\linewidth-20mm}
    {\fontsize{20}{24}\selectfont\bfseries\color{titolo} #1}\par\vspace{1mm}
    {\footnotesize\bfseries Indirizzo: IP14 -- MANUTENZIONE E ASSISTENZA TECNICA \quad -- \quad Classe 3MAT}
  \end{minipage}\par\vspace{2mm}
  {\color{titolo}\rule{\linewidth}{0.8pt}}\par\vspace{3mm}
  \begin{tcolorbox}[colback=grigio, colframe=grigio, arc=0pt, left=6pt, top=4pt, bottom=4pt]
    {\Large\bfseries #2}
  \end{tcolorbox}}

\newcommand{\nota}[1]{\begin{tcolorbox}[colback=grigio!45, colframe=grigio, arc=0pt, left=6pt, top=3pt, bottom=3pt]\small #1\end{tcolorbox}}
\newcommand{\tipo}[1]{\par\vspace{2mm}{\large\color{titolo} #1}\par\vspace{-1mm}{\color{titolo!60}\rule{\linewidth}{0.4pt}}\par\vspace{1mm}}
\newcommand{\esercizio}[2]{\par\vspace{3mm}\textbf{Esercizio #1.} \textbf{#2}\par}

% alimentazione: barra con etichetta
\newcommand{\rail}[2]{\draw[line width=1.2pt] ($#1+(-0.6,0)$) -- ($#1+(0.6,0)$); \node[above] at #1 {#2};}

\pagestyle{fancy}
\fancyhf{}
\renewcommand{\headrulewidth}{0pt}
\newcommand{\piede}[1]{\fancyfoot[L]{\scriptsize I.S.I.S. Sacile-Brugnera | sede: ``Piero Della Valentina'' | #1}
  \fancyfoot[R]{\scriptsize Pagina \thepage\ di \pageref{LastPage}}}

% Schemi degli esercizi (usati sia nell'esercitazione sia nelle soluzioni)
\newcommand{\partitoreV}[3]{% #1 = V, #2 = R1, #3 = R2
\begin{circuitikz}[scale=0.8, transform shape, line width=0.6pt]
  \rail{(0,4.4)}{$V = #1$}
  \draw (0,4.4) -- (0,4.2) to[R, l_={$R_1 = #2$}] (0,2.2) to[short, *-] (0,2.2)
        to[R, l_={$R_2 = #3$}] (0,0.2) -- (0,0) node[sground]{};
  \draw (0,2.2) to[short, -o] (1.4,2.2) node[right]{X};
\end{circuitikz}}

\newcommand{\partitoreI}[3]{% #1 = I, #2 = R1, #3 = R2
\begin{circuitikz}[scale=0.8, transform shape, line width=0.6pt]
  \draw (0,4.4) node[ocirc]{} to[short, i>^={$I = #1$}] (0,3.2);
  \draw (-1.6,3.2) -- (1.6,3.2); \draw (0,3.2) node[circ]{};
  \draw (-1.6,3.2) to[short, i>^={$I_1$}] (-1.6,2.6) to[R, l_={$R_1 = #2$}] (-1.6,0.6) -- (-1.6,0.3) -- (1.6,0.3);
  \draw (1.6,3.2) to[short, i>_={$I_2$}] (1.6,2.6) to[R, l^={$R_2 = #3$}] (1.6,0.6) -- (1.6,0.3);
  \draw (0,0.3) node[circ]{} -- (0,0.1) node[sground]{};
\end{circuitikz}}

\newcommand{\partitorePot}{%
\begin{circuitikz}[scale=0.8, transform shape, line width=0.6pt]
  \rail{(0,4.4)}{$V = 10\,\mathrm{V}$}
  \draw (0,4.4) -- (0,4.2) to[R, l_={$R_1 = 1\,\mathrm{k\Omega}$}] (0,2.2) to[short, *-] (0,2.2)
        to[vR, l_={$R_2 = 0 \div 4{,}7\,\mathrm{k\Omega}$}] (0,0.2) -- (0,0) node[sground]{};
  \draw (0,2.2) to[short, -o] (1.4,2.2) node[right]{X};
\end{circuitikz}}

\newcommand{\resistenzaSola}[2]{% #1 = V, #2 = R
\begin{circuitikz}[scale=0.8, transform shape, line width=0.6pt]
  \rail{(0,3.4)}{$V = #1$}
  \draw (0,3.4) -- (0,3.2) to[R, l_={$R = #2$}, i>^={$I$}] (0,0.8) -- (0,0.6) node[sground]{};
\end{circuitikz}}

\piede{Soluzioni esercitazione TEEA 3MAT}
\newcommand{\pas}[1]{\par\hspace*{4mm}$\displaystyle #1$\par}
\newcommand{\ris}[1]{\par\vspace{1mm}\hspace*{4mm}\colorbox{grigio}{\strut\ #1\ }\par}

\begin{document}

\intestazione{SOLUZIONI ESERCITAZIONE}{Partitori, potenza ed energia -- Soluzioni}

\nota{I valori sono arrotondati a tre cifre significative. Le resistenze in k$\Omega$ con le correnti in mA
danno direttamente le tensioni in V (V = k$\Omega$ $\cdot$ mA).}

\tipo{Partitore di tensione}

\esercizio{1}{$V = 24\U{V}$, $R_1 = 3{,}3\U{k\Omega}$, $R_2 = 4{,}7\U{k\Omega}$}
\pas{I = \frac{V}{R_1 + R_2} = \frac{24\U{V}}{3{,}3 + 4{,}7\U{k\Omega}} = \frac{24}{8} = 3\U{mA}}
\pas{V_X = V \cdot \frac{R_2}{R_1 + R_2} = 24 \cdot \frac{4{,}7}{8} = 14{,}1\U{V} \qquad V_1 = 24 \cdot \frac{3{,}3}{8} = 9{,}9\U{V} \qquad 9{,}9 + 14{,}1 = 24\U{V}\ \checkmark}
\pas{P_1 = V_1 \cdot I = 9{,}9 \cdot 0{,}003 = 29{,}7\U{mW} \qquad P_2 = V_X \cdot I = 14{,}1 \cdot 0{,}003 = 42{,}3\U{mW}}
\pas{P_{tot} = V \cdot I = 24 \cdot 0{,}003 = 72\U{mW} = 29{,}7 + 42{,}3\ \checkmark}
\ris{$I = 3\U{mA}$ \quad $V_1 = 9{,}9\U{V}$ \quad $V_X = 14{,}1\U{V}$ \quad $P_1 = 29{,}7\U{mW}$ \quad $P_2 = 42{,}3\U{mW}$ \quad $P_{tot} = 72\U{mW}$}

\esercizio{2}{$V = 15\U{V}$, $V_X = 5\U{V}$, $R_2 = 2{,}2\U{k\Omega}$}
\pas{\text{Su } R_2 \text{ ci sono } 5\U{V}: \ I = \frac{V_X}{R_2} = \frac{5}{2{,}2} = 2{,}27\U{mA} \qquad \text{su } R_1: \ V_1 = 15 - 5 = 10\U{V}}
\pas{R_1 = \frac{V_1}{I} = \frac{10\U{V}}{2{,}27\U{mA}} = 4{,}4\U{k\Omega} \qquad \Big(\text{oppure } R_1 = R_2 \cdot \frac{V - V_X}{V_X} = 2{,}2 \cdot \frac{10}{5} = 4{,}4\U{k\Omega}\Big)}
\pas{\text{valore E12 più vicino: } 4{,}7\U{k\Omega} \quad\Rightarrow\quad V_X = 15 \cdot \frac{2{,}2}{4{,}7 + 2{,}2} = 15 \cdot \frac{2{,}2}{6{,}9} = 4{,}78\U{V}}
\ris{$R_1 = 4{,}4\U{k\Omega}$ \quad commerciale $4{,}7\U{k\Omega}$ \quad $V_X = 4{,}78\U{V}$}

\tipo{Partitore di corrente}

\esercizio{3}{$I = 30\U{mA}$, $R_1 = 2\U{k\Omega}$, $R_2 = 3\U{k\Omega}$}
\pas{I_1 = I \cdot \frac{R_2}{R_1 + R_2} = 30 \cdot \frac{3}{5} = 18\U{mA} \qquad I_2 = I \cdot \frac{R_1}{R_1 + R_2} = 30 \cdot \frac{2}{5} = 12\U{mA} \qquad 18 + 12 = 30\ \checkmark}
\pas{R_{eq} = \frac{R_1 \cdot R_2}{R_1 + R_2} = \frac{2 \cdot 3}{5} = 1{,}2\U{k\Omega} \qquad V = R_{eq} \cdot I = 1{,}2\U{k\Omega} \cdot 30\U{mA} = 36\U{V}}
\pas{\text{controllo: } R_1 \cdot I_1 = 2 \cdot 18 = 36\U{V} \qquad R_2 \cdot I_2 = 3 \cdot 12 = 36\U{V}\ \checkmark}
\ris{$I_1 = 18\U{mA}$ \quad $I_2 = 12\U{mA}$ \quad $R_{eq} = 1{,}2\U{k\Omega}$ \quad $V = 36\U{V}$}

\esercizio{4}{$I = 50\U{mA}$, $R_1 = 100\U{\Omega}$, $R_2 = 400\U{\Omega}$}
\pas{I_1 = 50 \cdot \frac{400}{500} = 40\U{mA} \qquad I_2 = 50 \cdot \frac{100}{500} = 10\U{mA} \qquad V = R_1 \cdot I_1 = 100 \cdot 0{,}04 = 4\U{V}}
\pas{P_1 = V \cdot I_1 = 4 \cdot 0{,}04 = 0{,}16\U{W} \qquad P_2 = V \cdot I_2 = 4 \cdot 0{,}01 = 0{,}04\U{W}}
\pas{\text{con margine: } 2 \cdot 0{,}16 = 0{,}32\U{W} > 0{,}25\U{W} \ \Rightarrow\ \text{per } R_1 \text{ serve un resistore da } 1/2\U{W}}
\ris{$I_1 = 40\U{mA}$ \quad $I_2 = 10\U{mA}$ \quad $V = 4\U{V}$ \quad $P_1 = 0{,}16\U{W}$ \quad $P_2 = 0{,}04\U{W}$ \quad $R_1$: meglio $1/2\U{W}$}

\newpage
\tipo{Potenziometro}

\esercizio{5}{$V = 10\U{V}$, $R_1 = 1\U{k\Omega}$ fissa, $R_2$ potenziometro da $0$ a $4{,}7\U{k\Omega}$}
\pas{\text{a) } R_2 = 0: \ V_{X\,min} = 0\U{V}, \quad I_{max} = \frac{V}{R_1} = \frac{10}{1} = 10\U{mA} \qquad R_2 = 4{,}7\U{k\Omega}: \ V_{X\,max} = 10 \cdot \frac{4{,}7}{5{,}7} = 8{,}25\U{V}}
\pas{\text{b) per arrivare fino a } V = 10\U{V} \text{ il potenziometro va al posto di } R_1 \text{ (verso l'alimentazione), con } R_2 \text{ fissa}}
\pas{V_{X\,min} = V \cdot \frac{R_2}{R_{1\,max} + R_2} = 2\U{V} \ \Rightarrow\ R_{1\,max} + 1 = \frac{10 \cdot 1}{2} = 5 \ \Rightarrow\ R_{1\,max} = 4\U{k\Omega}}
\pas{\text{con un potenziometro commerciale da } 4{,}7\U{k\Omega}: \ V_{X\,min} = 10 \cdot \frac{1}{5{,}7} = 1{,}75\U{V}}
\ris{a) $0 \div 8{,}25\U{V}$, \ $I_{max} = 10\U{mA}$ \qquad b) al posto di $R_1$, \ $R_{1\,max} = 4\U{k\Omega}$ (commerciale $4{,}7\U{k\Omega}$)}

\tipo{Potenza ed energia}

\esercizio{6}{$R = 470\U{\Omega}$, $V = 12\U{V}$}
\pas{I = \frac{V}{R} = \frac{12}{470} = 0{,}0255\U{A} = 25{,}5\U{mA} \qquad P = \frac{V^2}{R} = \frac{144}{470} = 0{,}306\U{W}}
\pas{0{,}306\U{W} > 0{,}25\U{W} \Rightarrow 1/4\U{W} \text{ non va bene;} \quad 2 \cdot 0{,}306 = 0{,}61\U{W} > 0{,}5\U{W} \Rightarrow \text{ si sceglie } 1\U{W}}
\ris{$I = 25{,}5\U{mA}$ \quad $P = 0{,}306\U{W}$ \quad $1/4\U{W}$ no \quad scelta: $1\U{W}$}

\esercizio{7}{$R = 100\U{\Omega}$, $I = 0{,}2\U{A}$, $t = 5\U{h}$}
\pas{V = R \cdot I = 100 \cdot 0{,}2 = 20\U{V} \qquad P = R \cdot I^2 = 100 \cdot 0{,}2^2 = 4\U{W} \quad (= V \cdot I = 20 \cdot 0{,}2)}
\pas{E = P \cdot t = 4\U{W} \cdot 5\U{h} = 20\U{Wh} = 0{,}02\U{kWh}}
\ris{$V = 20\U{V}$ \quad $P = 4\U{W}$ \quad $E = 20\U{Wh} = 0{,}02\U{kWh}$}

\esercizio{8}{$P = 1500\U{W}$, $V = 230\U{V}$, $4\U{h}$ al giorno per $30$ giorni, $0{,}25$ euro/kWh}
\pas{I = \frac{P}{V} = \frac{1500}{230} = 6{,}52\U{A} \qquad R = \frac{V^2}{P} = \frac{230^2}{1500} = \frac{52\,900}{1500} = 35{,}3\U{\Omega}}
\pas{t = 4 \cdot 30 = 120\U{h} \qquad E = P \cdot t = 1500 \cdot 120 = 180\,000\U{Wh} = 180\U{kWh}}
\pas{\text{costo} = 180\U{kWh} \cdot 0{,}25 \text{ euro/kWh} = 45 \text{ euro}}
\ris{$I = 6{,}52\U{A}$ \quad $R = 35{,}3\U{\Omega}$ \quad $E = 180\,000\U{Wh} = 180\U{kWh}$ \quad costo $= 45$ euro}

\end{document}
